\documentclass[11pt,aspectratio=169]{beamer}

%%%%%%%%%%%%%%%%%%%%%%%%%%%%%%
% Session 4: Wording, Response Options, Order, and Sensitive Questions
% LBJ Survey Design and Analysis, Fall 2026
%%%%%%%%%%%%%%%%%%%%%%%%%%%%%%

\usetheme[titleformat=smallcaps,progressbar=frametitle,block=fill]{metropolis}

\usepackage[utf8]{inputenc}
\usepackage[T1]{fontenc}
\usepackage{booktabs}
\usepackage{tabularx}
\usepackage{array}
\usepackage{multirow}
\usepackage{ragged2e}
\usepackage{tikz}
\usetikzlibrary{arrows.meta,positioning,calc,fit}
\usepackage{amsmath}
\usepackage{amssymb}
\usepackage{hyperref}
\usepackage{xcolor}

%%%%%%%%%%%%%%%%%%%%%%%%%%%%%%
% Course colors
%%%%%%%%%%%%%%%%%%%%%%%%%%%%%%
\xdefinecolor{custom_blue}{rgb}{0.01, 0.31, 0.52}
\xdefinecolor{custom_darkBlue}{rgb}{0.20, 0.20, 0.39}
\xdefinecolor{custom_orange}{rgb}{0.96, 0.57, 0.42}
\xdefinecolor{custom_green}{rgb}{0, 0.47, 0.52}
\xdefinecolor{custom_red}{rgb}{0.58, 0.32, 0.32}
\xdefinecolor{custom_lightGray}{rgb}{0.78, 0.80, 0.80}
\xdefinecolor{custom_darkGray}{rgb}{0.35, 0.39, 0.43}

\setbeamercolor{alerted text}{fg=custom_orange}
\setbeamercolor{item}{fg=custom_blue}
\setbeamercolor{subitem}{fg=custom_blue}
\setbeamerfont{frametitle}{size=\small}
\setbeamerfont{title}{size=\Large}
\setbeamerfont{subtitle}{size=\normalsize}
\setbeamertemplate{navigation symbols}{}
\setbeamertemplate{footline}{%
  \raisebox{4pt}{\makebox[\paperwidth][r]{\makebox[24pt]{\color{gray}\scriptsize\insertframenumber}}}
}

\newcommand{\hl}[1]{\textcolor{custom_blue}{#1}}
\newcommand{\orange}[1]{\textcolor{custom_orange}{#1}}
\newcommand{\green}[1]{\textcolor{custom_green}{#1}}
\newcommand{\red}[1]{\textcolor{custom_red}{#1}}
\newcommand{\gray}[1]{\textcolor{custom_darkGray}{#1}}
\newcommand{\ct}[1]{\vfill{\tiny\color{custom_darkGray}#1}}

\newcommand{\yourturn}[1]{%
\setbeamercolor{block title}{fg=white,bg=custom_blue}
\setbeamercolor{block body}{fg=black,bg=custom_blue!10}
\begin{block}{Your turn}
#1
\end{block}}

\newcommand{\groupwork}[2]{%
\setbeamercolor{block title}{fg=white,bg=custom_green}
\setbeamercolor{block body}{fg=black,bg=custom_green!10}
\begin{block}{Small-group work: #1}
#2
\end{block}}

\newcommand{\keypoint}[2]{%
\setbeamercolor{block title}{fg=white,bg=custom_orange}
\setbeamercolor{block body}{fg=black,bg=custom_orange!10}
\begin{block}{#1}
#2
\end{block}}

\title{Wording, Response Options, Order, and Sensitive Questions}
\subtitle{From a measurement instrument to a respondent's answer}
\author{Sergio I. Garc\'ia-R\'ios}
\institute{LBJ School of Public Affairs \;|\; Survey Design and Analysis}
\date{September 21, 2026}

\begin{document}

% ------------------------------------------------------------------
\maketitle

% ------------------------------------------------------------------
\begin{frame}[standout]
\Large Last week: What are we trying to measure?\\[0.8em]
\Large Today: What happens when a \alert{human respondent} encounters our measure?
\end{frame}

% ------------------------------------------------------------------
\begin{frame}{Where we stopped}
Last week we moved from:

\begin{center}
\large Research objective $\rightarrow$ construct $\rightarrow$ dimensions $\rightarrow$ indicators
\end{center}

\pause
Today we finish the right-hand side of the chain:

\begin{center}
\large indicators $\rightarrow$ \alert{question} $\rightarrow$ \alert{respondent process} $\rightarrow$ observed response
\end{center}

\vfill
\keypoint{Continuity, not recap}{We are not starting a second topic. We are opening up the mechanism that turns a question into data.}
\end{frame}

% ------------------------------------------------------------------
\begin{frame}{Finish the reading workshop}
\small
\begin{columns}[T]
\column{0.32\textwidth}
\textbf{Group 1: Groves Ch. 7}

What model of survey measurement is being proposed? Where can error enter?

\column{0.32\textwidth}
\textbf{Group 2: Krosnick \& Presser}

What principles govern the move from a construct to a question and response format?

\column{0.32\textwidth}
\textbf{Group 3: Pew}

How are those principles translated into professional questionnaire practice?
\end{columns}

\vspace{1em}
\groupwork{8 minutes}{Return to the reading with your group. Your job is to teach us the \textbf{argument}, not summarize pages.}
\end{frame}

% ------------------------------------------------------------------
\begin{frame}{Four questions for every reading}
Each group should be ready to answer:

\begin{enumerate}
\item What is the central \alert{methodological claim}?
\item What assumption about respondents or measurement is underneath it?
\item What \alert{design decision} follows from that claim?
\item What kind of error could result if we ignore it?
\end{enumerate}

\vfill
\yourturn{When another group reports, prepare one challenge: \textit{When might that recommendation fail, or create a different problem?}}
\end{frame}

% ------------------------------------------------------------------
\begin{frame}{Cross-examination}
Graduate methods should not end with ``best practices.''

\begin{itemize}
\item Does the recommendation depend on the \alert{construct}?
\item Does it depend on the \alert{population}?
\item Does it depend on \alert{mode}?
\item What is the tradeoff between respondent burden and measurement precision?
\item Could reducing one source of error increase another?
\end{itemize}

\vfill
\keypoint{The standard}{A design rule is useful only if you understand the mechanism it is trying to control.}
\end{frame}

% ------------------------------------------------------------------
\begin{frame}{What Groves adds to the measurement chain}
\centering
\begin{tikzpicture}[
  node distance=0.48cm,
  box/.style={draw=custom_blue, rounded corners, minimum height=.95cm, minimum width=2.25cm, align=center, fill=custom_blue!6, font=\scriptsize},
  arr/.style={-{Latex[length=2.5mm]}, thick, draw=custom_darkGray}
]
\node[box] (q) {Question presented};
\node[box, right=of q] (comp) {Comprehend};
\node[box, right=of comp] (ret) {Retrieve};
\node[box, right=of ret] (jud) {Judge / estimate};
\node[box, right=of jud] (map) {Map to response};
\draw[arr] (q) -- (comp);
\draw[arr] (comp) -- (ret);
\draw[arr] (ret) -- (jud);
\draw[arr] (jud) -- (map);
\end{tikzpicture}

\vspace{1.4em}
\pause
Each stage can introduce error --- even if the respondent is cooperative and trying to answer well.

\vfill
\ct{Groves et al. (2009), Ch. 7; Tourangeau, Rips, and Rasinski (2000).}
\end{frame}

% ------------------------------------------------------------------
\begin{frame}{A perfect sample does not solve measurement error}
Suppose we draw a flawless probability sample of the target population.

\pause
\begin{center}
\Large $\widehat{Y}$ can still be a very precise estimate\\[0.4em]
\Large of the \alert{wrong quantity}.
\end{center}

\pause
Why?
\begin{itemize}
\item the construct was underspecified;
\item the item captures only one dimension;
\item respondents interpret the question differently;
\item retrieval is difficult;
\item the response options distort mapping;
\item context changes the judgment being reported.
\end{itemize}
\end{frame}

% ------------------------------------------------------------------
\begin{frame}{Open and closed questions: same construct?}
\small
\begin{columns}[T]
\column{0.48\textwidth}
\textbf{Open}

\medskip
What do you think is the most important problem facing the country today?

\column{0.48\textwidth}
\textbf{Closed}

\medskip
Which of the following is the most important problem facing the country today?

\begin{itemize}
\item Cost of living
\item Immigration
\item Health care
\item Crime
\item Education
\item Something else
\end{itemize}
\end{columns}

\vfill
\yourturn{Are these two instruments measuring the same thing? Defend your answer.}
\end{frame}

% ------------------------------------------------------------------
\begin{frame}{Open versus closed is a measurement choice}
\small
\begin{tabularx}{\textwidth}{>{\bfseries}lXX}
\toprule
 & Open-ended & Closed-ended \\
\midrule
Cognitive task & Recall / generate & Recognize / select \\
Researcher control & Low & High \\
Respondent burden & Often higher & Often lower \\
Standardization & Lower & Higher \\
Risk & Coding + articulation & Omission + framing \\
What becomes salient? & Respondent-generated & Researcher-provided \\
\bottomrule
\end{tabularx}

\vfill
\keypoint{Not a commandment}{Choose the format that best matches the inferential task and the respondent's cognitive task.}
\ct{Krosnick and Presser (2010); Pew Research Center, ``Writing Survey Questions.''}
\end{frame}

% ------------------------------------------------------------------
\begin{frame}{Response options are part of the measure}
Imagine asking:

\begin{center}
\large How often have you felt very irritated during the past month?
\end{center}

Two researchers provide very different frequency ranges.

\pause
\begin{center}
\Large Why might the \alert{range itself} alter the answer?
\end{center}

\vfill
\keypoint{Bridge to Schwarz}{Respondents can use response alternatives as information about what the researcher means, what counts as typical, and which distinctions matter.}
\end{frame}

% ------------------------------------------------------------------
\begin{frame}[standout]
\Large The questionnaire is not a neutral container.\\[0.8em]
\Large It is part of the \alert{stimulus} respondents are reacting to.
\end{frame}

% ------------------------------------------------------------------
\section{Schwarz: how questions shape answers}

% ------------------------------------------------------------------
\begin{frame}{This week's readings}
\begin{enumerate}
\item \textbf{Schwarz (1999)} --- cognitive and communicative processes underlying self-reports.
\item \textbf{Pew Methods 101 (2018)} --- practical consequences for wording and ordering.
\item \textbf{Tourangeau \& Yan (2007)} --- what changes when respondents may not want to report truthfully.
\end{enumerate}

\vfill
\keypoint{The common problem}{The observed answer depends not only on the respondent's underlying state, but on the task we create for them.}
\end{frame}

% ------------------------------------------------------------------
\begin{frame}{Schwarz's central claim}
\large
Self-reports of attitudes and behaviors are shaped by features of the research instrument --- including \alert{wording, format, and context}.

\vspace{1em}
\pause
The important part is not merely that ``wording matters.''

\begin{center}
\Large The effects arise from \alert{systematic cognitive and communicative processes}.
\end{center}

\vfill
\ct{Schwarz (1999), pp. 93--105.}
\end{frame}

% ------------------------------------------------------------------
\begin{frame}{Respondents answer more than the literal sentence}
Consider:

\begin{center}
\large ``How often do you exercise regularly?''
\end{center}

\pause
The respondent must infer:
\begin{itemize}
\item What counts as \alert{exercise}?
\item What does \alert{regularly} mean?
\item Which time period is relevant?
\item What level of precision is expected?
\item Why is the researcher asking this question here?
\end{itemize}

\vfill
\keypoint{Literal + pragmatic meaning}{Survey answering is a constrained conversation. Respondents infer the researcher's intended meaning from the words and the surrounding instrument.}
\end{frame}

% ------------------------------------------------------------------
\begin{frame}{The questionnaire as a conversation}
Schwarz draws on the logic of ordinary conversation: respondents generally assume that the researcher is trying to be meaningful and informative.

\pause
That means features we may treat as ``formatting'' can be interpreted as \alert{information}:

\begin{itemize}
\item response categories;
\item numeric endpoints;
\item examples included in the stem;
\item reference periods;
\item questions that came immediately before;
\item what the questionnaire appears to be about.
\end{itemize}

\vfill
\yourturn{What information might a respondent infer from a scale that ranges from ``0 times'' to ``more than 10 times''?}
\end{frame}

% ------------------------------------------------------------------
\begin{frame}{Behavior reports: retrieval or estimation?}
Suppose I ask:

\begin{center}
\large ``How many times have you eaten at a restaurant in the past six months?''
\end{center}

\pause
Are you retrieving six months of discrete episodes from memory?

\pause
Probably not. For frequent, mundane behaviors, respondents often construct an estimate using:

\begin{itemize}
\item typical rates (``about twice a week'');
\item partial retrieval;
\item category-based inference;
\item cues provided by the question and response options.
\end{itemize}

\vfill
\ct{Schwarz (1999), section on behavioral frequency reports.}
\end{frame}

% ------------------------------------------------------------------
\begin{frame}{Reference periods change the cognitive task}
Compare:

\begin{enumerate}
\item How many times did you eat at a restaurant \alert{during the past seven days}?
\item How many times did you eat at a restaurant \alert{during the past six months}?
\end{enumerate}

\pause
\begin{center}
\Large Same behavior. Different retrieval problem.
\end{center}

\vfill
\keypoint{Design implication}{A reference period is not just a calendar boundary. It determines whether respondents can retrieve events or must estimate them.}
\end{frame}

% ------------------------------------------------------------------
\begin{frame}{Small-group Schwarz lab}
\small
\begin{columns}[T]
\column{0.32\textwidth}
\textbf{Group A: Comprehension}

``How often have you exercised regularly during the past year?''

Diagnose literal and pragmatic ambiguity.

\column{0.32\textwidth}
\textbf{Group B: Retrieval}

Compare a 7-day and 6-month restaurant-use question.

Explain how the response strategy changes.

\column{0.32\textwidth}
\textbf{Group C: Response scale}

Design two different frequency scales for the same behavior.

Predict how each scale could alter reports.
\end{columns}

\vspace{1em}
\groupwork{10 minutes}{Do not merely rewrite the question. Identify the \textbf{cognitive or communicative mechanism} that creates the problem.}
\end{frame}

% ------------------------------------------------------------------
\begin{frame}{Schwarz lab: debrief standard}
For each diagnosis, we want four pieces:

\begin{enumerate}
\item What is the respondent trying to infer or retrieve?
\item What feature of the instrument changes that task?
\item What direction of error or heterogeneity might result?
\item What redesign would reduce the problem --- and what might that redesign sacrifice?
\end{enumerate}

\vfill
\keypoint{Graduate-level diagnosis}{``Bad wording'' is not a mechanism. Tell me \textbf{why} the observed answer changes.}
\end{frame}

% ------------------------------------------------------------------
\begin{frame}{Response scales can communicate meaning}
A response scale does at least two jobs:

\begin{enumerate}
\item It gives respondents a way to \alert{map} a judgment into an answer.
\item It signals what the researcher thinks is a plausible or meaningful range.
\end{enumerate}

\pause
Therefore, two scales with the same verbal question can induce different interpretations of:

\begin{itemize}
\item what counts as ``often'';
\item what counts as ``normal'';
\item which distinctions the researcher wants the respondent to make.
\end{itemize}

\ct{Schwarz (1999).}
\end{frame}

% ------------------------------------------------------------------
\begin{frame}{Question order can change the judgment itself}
Suppose respondents answer:

\begin{enumerate}
\item How satisfied are you with your dating life?
\item How satisfied are you with your life as a whole?
\end{enumerate}

versus the reverse order.

\pause
The first question can make dating information unusually accessible when answering the second.

\pause
\begin{center}
\large The later answer may reflect a different \alert{information set}, not a different person.
\end{center}

\vfill
\ct{Schwarz (1999), discussion of context effects in attitude measurement.}
\end{frame}

% ------------------------------------------------------------------
\begin{frame}{Assimilation and contrast}
Context effects are not always ``more of the same.''

\begin{columns}[T]
\column{0.48\textwidth}
\textbf{Assimilation}

Previously activated information is \alert{included} in the target judgment.

\medskip
Example logic: ``My dating life is good, so my life feels good.''

\column{0.48\textwidth}
\textbf{Contrast}

Previously activated information becomes a \alert{comparison standard} or is excluded from the target.

\medskip
Example logic: ``Compared with that unusually good case, this seems worse.''
\end{columns}

\vfill
\keypoint{Why order effects are substantive}{Question order can change what respondents treat as evidence for the judgment.}
\end{frame}

% ------------------------------------------------------------------
\begin{frame}{Order experiment: predict before we test}
Imagine a survey contains these two modules:

\begin{itemize}
\item detailed questions about personal financial strain;
\item overall evaluation of one's quality of life.
\end{itemize}

\yourturn{Would you expect the quality-of-life measure to differ if the financial-strain module comes first? Under what theory?}

\pause
Possible mechanisms:
\begin{itemize}
\item accessibility / priming;
\item conversational inference about what ``quality of life'' is supposed to mean;
\item assimilation or contrast;
\item fatigue or satisficing --- a different mechanism entirely.
\end{itemize}
\end{frame}

% ------------------------------------------------------------------
\begin{frame}[standout]
\Large Survey respondents do not simply ``reveal'' an answer.\\[0.8em]
\Large They \alert{construct} an answer under the conditions we design.
\end{frame}

% ------------------------------------------------------------------
\section{Pew: translating mechanisms into design decisions}

% ------------------------------------------------------------------
\begin{frame}{Pew turns theory into production rules}
Pew emphasizes that sound questionnaire design requires attention to:

\begin{itemize}
\item clear, specific, answerable wording;
\item exhaustive and mutually exclusive categories;
\item question order and context;
\item wording and context when measuring change over time;
\item mode when interpreting trends;
\item social desirability for sensitive behaviors and attitudes.
\end{itemize}

\vfill
\keypoint{Ask the deeper question}{What failure in the response process is each practice trying to prevent?}
\ct{Pew Research Center (2018), Methods 101: Survey Question Wording; Pew, ``Writing Survey Questions.''}
\end{frame}

% ------------------------------------------------------------------
\begin{frame}{Map the ``rule'' to the mechanism}
\small
\begin{tabularx}{\textwidth}{>{\bfseries}lX}
\toprule
Design practice & Measurement problem it is trying to control \\
\midrule
Clear, specific wording & Heterogeneous comprehension \\
Concrete reference period & Retrieval / estimation error \\
Exhaustive categories & Forced misclassification \\
Non-overlapping categories & Ambiguous response mapping \\
Attention to order & Context-dependent judgment \\
Same wording for trends & Changes in the measurement instrument \\
Self-administration for sensitive items & Social desirability / disclosure pressure \\
\bottomrule
\end{tabularx}
\end{frame}

% ------------------------------------------------------------------
\begin{frame}{A wording change can change the estimand}
Pew's historical example asks about support for U.S. military action in Iraq in 2003.

\medskip
One version asked about military action to end Saddam Hussein's rule. Another explicitly added the possibility of thousands of U.S. casualties.

\pause
\begin{center}
\Large Support: \alert{68\%} without the casualty clause\\[0.2em]
\Large versus \alert{43\%} with it.
\end{center}

\pause
\yourturn{Is one version ``biased''? Or are they asking respondents to evaluate different policy scenarios?}

\ct{Pew Research Center, ``Writing Survey Questions,'' historical January 2003 wording experiment.}
\end{frame}

% ------------------------------------------------------------------
\begin{frame}{Not every wording effect is a defect}
A wording difference can reveal several different things:

\begin{enumerate}
\item one version is genuinely confusing or leading;
\item the versions activate different considerations;
\item the versions define different hypothetical conditions;
\item the construct itself is conditional on information respondents receive.
\end{enumerate}

\pause
\keypoint{Analytic discipline}{Before calling a wording effect ``bias,'' specify the construct and the estimand each version is supposed to measure.}
\end{frame}

% ------------------------------------------------------------------
\begin{frame}{Trend measurement creates a different problem}
Suppose an item has been asked identically for 15 years, but today you believe you can write it better.

\pause
Two goals now conflict:

\begin{columns}[T]
\column{0.48\textwidth}
\textbf{Improve current measurement}
\begin{itemize}
\item clearer wording
\item better categories
\item new mode / device realities
\end{itemize}

\column{0.48\textwidth}
\textbf{Preserve comparability}
\begin{itemize}
\item same wording
\item same context
\item same response task
\end{itemize}
\end{columns}

\vfill
\keypoint{Measurement invariance problem}{A better item today may break the time series you care about.}
\end{frame}

% ------------------------------------------------------------------
\begin{frame}{Small-group design decision: fix it or preserve it?}
A client has a long-running question:

\begin{center}
\large ``How would you rate the state of your household finances today: excellent, good, only fair, or poor?''
\end{center}

The client wants to add explicit language about ``paying bills and handling unexpected expenses'' because they think respondents interpret ``finances'' inconsistently.

\groupwork{8 minutes}{Choose a recommendation. Keep the old item, replace it, run both, or do something else. State what evidence you would want before changing the trend.}
\end{frame}

% ------------------------------------------------------------------
\begin{frame}{Debrief: there is no free redesign}
A defensible transition strategy might involve:

\begin{itemize}
\item a split-ballot bridge experiment;
\item running old and new measures in parallel for one or more waves;
\item testing mode-by-wording interactions;
\item documenting a series break rather than disguising it;
\item deciding explicitly whether the priority is \alert{level estimation} or \alert{change over time}.
\end{itemize}

\vfill
\keypoint{Production lesson}{Questionnaire design includes preserving the meaning of a measure across versions of the instrument.}
\end{frame}

% ------------------------------------------------------------------
\section{Tourangeau and Yan: when respondents may edit the truth}

% ------------------------------------------------------------------
\begin{frame}{Sensitive questions are not one category}
Tourangeau and Yan distinguish several ways a question can be sensitive.

\begin{enumerate}
\item \alert{Intrusiveness}: the topic itself feels private or inappropriate to ask.
\item \alert{Threat of disclosure}: truthful reporting could have consequences if disclosed.
\item \alert{Social desirability}: one response is normatively more acceptable than another.
\end{enumerate}

\pause
\yourturn{Is household income sensitive for the same reason as illegal drug use? What about voting?}

\ct{Tourangeau and Yan (2007), pp. 859--870.}
\end{frame}

% ------------------------------------------------------------------
\begin{frame}{Sensitivity is relational and situational}
A question is not simply ``sensitive'' in the abstract.

Sensitivity can depend on:
\begin{itemize}
\item what the respondent's true answer is;
\item who appears to be asking;
\item whether an interviewer is present;
\item who else is nearby;
\item perceived confidentiality;
\item legal, social, or interpersonal consequences;
\item norms within the respondent's social context.
\end{itemize}

\vfill
\keypoint{Tourangeau \& Yan's important move}{Misreporting is often better understood as a \textbf{situational response strategy} than as a stable tendency of ``dishonest respondents.''}
\end{frame}

% ------------------------------------------------------------------
\begin{frame}{Sensitive questions add an editing stage}
\centering
\begin{tikzpicture}[
  node distance=0.38cm,
  box/.style={draw=custom_blue, rounded corners, minimum height=.88cm, minimum width=1.7cm, align=center, fill=custom_blue!6, font=\scriptsize},
  edit/.style={draw=custom_red, rounded corners, minimum height=.88cm, minimum width=1.7cm, align=center, fill=custom_red!8, font=\scriptsize},
  arr/.style={-{Latex[length=2.4mm]}, thick, draw=custom_darkGray}
]
\node[box] (comp) {Comprehend};
\node[box, right=of comp] (ret) {Retrieve};
\node[box, right=of ret] (jud) {Judge};
\node[edit, right=of jud] (edit) {Edit /\ self-present};
\node[box, right=of edit] (map) {Report};
\draw[arr] (comp) -- (ret);
\draw[arr] (ret) -- (jud);
\draw[arr] (jud) -- (edit);
\draw[arr] (edit) -- (map);
\end{tikzpicture}

\vspace{1.2em}
For sensitive items, respondents can know the truthful answer perfectly well and still choose not to report it.

\vfill
\ct{Tourangeau and Yan (2007), conclusion.}
\end{frame}

% ------------------------------------------------------------------
\begin{frame}{Sensitivity can produce different kinds of missingness and error}
\small
\begin{tabularx}{\textwidth}{>{\bfseries}lX}
\toprule
Response & What happened? \\
\midrule
Unit nonresponse & Person declines the survey, perhaps because of the topic. \\
Item nonresponse & Person participates but refuses / skips the sensitive item. \\
Misreporting & Person answers, but edits in a socially safer direction. \\
Heaping / categorization & Person provides a less precise answer that feels safer. \\
Mode selection / breakoff & Person changes engagement or leaves when privacy feels insufficient. \\
\bottomrule
\end{tabularx}

\vfill
\keypoint{Why this matters}{Item response is not evidence that measurement succeeded.}
\end{frame}

% ------------------------------------------------------------------
\begin{frame}{Directional error is especially dangerous}
Ordinary reporting errors can sometimes add noise in both directions.

Sensitive-question editing often has a directional component:

\begin{itemize}
\item socially desirable behavior may be \alert{overreported};
\item socially undesirable behavior may be \alert{underreported};
\item stigmatized attitudes may be softened;
\item threatening disclosures may be concealed.
\end{itemize}

\pause
\begin{center}
\large That creates \alert{bias in estimates}, not merely larger variance.
\end{center}

\ct{Tourangeau and Yan (2007).}
\end{frame}

% ------------------------------------------------------------------
\begin{frame}{Design features can change reporting}
The evidence reviewed by Tourangeau and Yan implies that reporting is responsive to survey design.

Common levers include:
\begin{itemize}
\item self-administration rather than interviewer administration;
\item physical / social privacy;
\item credible confidentiality protections;
\item response formats that reduce the interpersonal act of disclosure;
\item wording that lowers unnecessary judgment or threat;
\item specialized techniques for highly sensitive behaviors.
\end{itemize}

\vfill
\keypoint{Important caution}{No technique makes sensitivity disappear. Each intervention changes the response task and may introduce new costs.}
\end{frame}

% ------------------------------------------------------------------
\begin{frame}{Small-group sensitive-question lab}
Your client wants an accurate measure of household financial strain.

They want to ask about:
\begin{itemize}
\item annual household income;
\item missed housing or utility payments;
\item debt burden;
\item receipt of public assistance;
\item inability to pay for necessities.
\end{itemize}

\groupwork{12 minutes}{Design a short module that you would be willing to defend methodologically and ethically.}
\end{frame}

% ------------------------------------------------------------------
\begin{frame}{Sensitive-question lab: required decisions}
Your group must specify:

\begin{enumerate}
\item Which items are sensitive --- and \alert{why}?
\item Where do you expect item nonresponse versus misreporting?
\item Would interviewer presence change the problem?
\item Exact values or categories? Why?
\item What introduction, privacy cue, or response design would you use?
\item What information are you willing to sacrifice to reduce respondent burden or threat?
\item How would you know whether your redesign actually improved measurement?
\end{enumerate}
\end{frame}

% ------------------------------------------------------------------
\begin{frame}{Debrief: ``more disclosure'' is not the only objective}
A strong design has to balance at least four goals:

\begin{enumerate}
\item \alert{Validity}: does the response correspond to the target construct?
\item \alert{Accuracy}: are respondents reporting what actually happened?
\item \alert{Respondent protection}: are we minimizing unnecessary risk and intrusion?
\item \alert{Usability}: can the resulting data support the intended analysis?
\end{enumerate}

\vfill
\keypoint{Ethics and measurement are linked}{A design that makes respondents feel unsafe can be both ethically weaker and methodologically worse.}
\end{frame}

% ------------------------------------------------------------------
\begin{frame}{One integrated model}
\centering
\begin{tikzpicture}[
  node distance=0.28cm,
  box/.style={draw=custom_blue, rounded corners, minimum height=.88cm, minimum width=1.55cm, align=center, fill=custom_blue!6, font=\tiny},
  redbox/.style={draw=custom_red, rounded corners, minimum height=.88cm, minimum width=1.55cm, align=center, fill=custom_red!8, font=\tiny},
  arr/.style={-{Latex[length=2.3mm]}, thick, draw=custom_darkGray}
]
\node[box] (c) {Construct};
\node[box, right=of c] (q) {Question};
\node[box, right=of q] (co) {Comprehend};
\node[box, right=of co] (r) {Retrieve};
\node[box, right=of r] (j) {Judge};
\node[redbox, right=of j] (e) {Edit?};
\node[box, right=of e] (m) {Map / report};
\draw[arr] (c) -- (q);
\draw[arr] (q) -- (co);
\draw[arr] (co) -- (r);
\draw[arr] (r) -- (j);
\draw[arr] (j) -- (e);
\draw[arr] (e) -- (m);
\end{tikzpicture}

\vspace{1.2em}
\begin{itemize}
\item \textbf{Groves:} distinguish constructs, measurement, response, and error.
\item \textbf{Schwarz:} the instrument changes comprehension, retrieval, judgment, and mapping.
\item \textbf{Tourangeau \& Yan:} sensitivity can add motivated editing before reporting.
\item \textbf{Pew:} professional design choices try to control these mechanisms in practice.
\end{itemize}
\end{frame}

% ------------------------------------------------------------------
\begin{frame}{From an item to a defensible module}
For your own survey work, every proposed item should survive five questions:

\begin{enumerate}
\item What construct / dimension is this intended to measure?
\item What cognitive task does the respondent need to perform?
\item What information does the wording, scale, and context communicate?
\item Is there a reason the respondent might edit or withhold the answer?
\item What claim will we make from the resulting variable?
\end{enumerate}

\vfill
\keypoint{The production standard}{If you cannot answer these five questions, the item is not ready for fielding.}
\end{frame}

% ------------------------------------------------------------------
\begin{frame}{What we intentionally did not reduce this to}
\begin{center}
\Large ``Avoid double-barreled questions.''\\[0.5em]
\Large ``Use neutral wording.''\\[0.5em]
\Large ``Put sensitive questions at the end.''
\end{center}

\pause
Those can be useful heuristics. But the methodological goal is to understand:

\begin{center}
\large \alert{what response process the design creates, for whom, and with what consequences.}
\end{center}
\end{frame}

% ------------------------------------------------------------------
\begin{frame}{Four ideas to leave with}
\begin{enumerate}
\item A response is \alert{constructed}, not simply retrieved from a mental file.
\item Wording, response alternatives, and order can change the \alert{meaning of the task}.
\item Sensitive questions create not only nonresponse but \alert{motivated reporting error}.
\item Good questionnaire design is the design of the \alert{conditions under which measurement occurs}.
\end{enumerate}
\end{frame}

% ------------------------------------------------------------------
\begin{frame}{Exit ticket}
\large
You are handed a survey item that seems clear and grammatically clean.

\vspace{1em}
\yourturn{What are the \textbf{three most important questions} you would now ask before deciding that it is a good measure?}

\vfill
\small Next: cognitive testing and multilingual equivalence --- how we find out whether respondents are actually doing what we think they are doing.
\end{frame}

% ------------------------------------------------------------------
\begin{frame}{References for today's class}
\footnotesize
\begin{itemize}
\item Groves, Robert M., et al. 2009. \textit{Survey Methodology}, 2nd ed., Ch. 7.
\item Krosnick, Jon A., and Stanley Presser. 2010. ``Question and Questionnaire Design.'' In \textit{Handbook of Survey Research}, 2nd ed.
\item Pew Research Center. 2018. ``Methods 101: Survey Question Wording.''
\item Pew Research Center. ``Writing Survey Questions.''
\item Schwarz, Norbert. 1999. ``Self-Reports: How the Questions Shape the Answers.'' \textit{American Psychologist} 54(2): 93--105.
\item Tourangeau, Roger, and Ting Yan. 2007. ``Sensitive Questions in Surveys.'' \textit{Psychological Bulletin} 133(5): 859--883.
\end{itemize}
\end{frame}

\end{document}
